\section{Measuring a number you can defend}
\label{sec:method}

\noindent\emph{Status: draft -- this section is an argument, not a result. Nothing in it is computed. It is the
rule that the measurements of this report are held to, written after breaking that rule at least once each.
Every number used as an example is the report's own, and the section it came from is cited so that the reader
can check the break rather than take my word for the correction.}

Every number in this report arrived by the same three-step route. A field was stepped, a record was written to
disk, and somebody applied an arithmetic operation to that record and called the result a measurement. The
arithmetic is not the problem: dividing two fluxes, taking a logarithm, fitting an exponential, and inverting a
$2\times 2$ system are all legitimate operations, and all of them appear in what follows. The problem is the
noun. \emph{Measurement} carries an implicit claim of independence, and most of the numbers in this report
cannot support that claim: they are derived quantities, and a derived quantity inherits every assumption of the
derivation while looking, in a table, exactly like something that was observed. What follows is the audit of
that word, in the order in which the mistakes were made. It is written in the first person because the errors
are mine, and it is written as a list of rules because a rule can be seen to be broken, while a good intention
cannot.

\subsection{A number without provenance is not a number}
\label{sec:method-provenance}

Before the rules about derivation, there is a rule about the record itself. A number in this report is a
statement about a particular run of a particular mesh processed by a particular estimator, and unless all three
are named the number is not reproducible even in principle. Two runs of the same structure at the same mesh
differ in their truncation of the field outside the computational window, in the frequency they were driven at,
and in the length of the record that was transformed; two estimators applied to one run differ far more than
that, as the next subsection shows. A single printed value is therefore a summary of a triple, and a value whose
triple has been lost is a rumour.

This is why Section~\ref{sec:worked} is built around a four-column table whose middle two columns are the
provenance: the cell size and the mesh counts, the estimator, the record length and the drive frequency. It reads
as bureaucracy and it is the only part of that table which cannot be argued with. Where the provenance is
missing, the number is marked as such rather than reconstructed, because reconstructing it means choosing the
run that makes it most plausible, which is a way of fitting the provenance to the answer.

The most useful thing provenance does is to make a class of disagreement impossible. When two numbers in the
record do not match -- and several do not -- there are two separate questions: could these two numbers have come
from the same run, and do they in fact agree. These are routinely conflated, and I conflated them myself. A
published pairing of $|S_{11}|$ and $|S_{21}|$ once recorded a $7.3$ per cent passivity deficit, while the single
record that now carries both quantities gives $|s_{21}|^2 + |\gamma|^2 = 0.96228$, a $3.77$ per cent deficit at
the same frequency. My first note on the discrepancy claimed that the pairing \emph{could not} have come from a
single run. That claim was too strong and it was wrong: a single run does yield both numbers at the resonance,
and they are simply different numbers. The failure is a numeric mismatch, which is a fact about the record, and
not a provenance impossibility, which would have been a fact about the archive.

The distinction matters because the two failures have different remedies. A provenance impossibility is closed
by finding the missing run; a numeric mismatch is closed only by re-measuring, and until then it is reported as
an open discrepancy. Saying the strong thing when the weak thing is true is not caution. It is a second error,
committed in the direction of looking rigorous.

\subsection{A fit is not a measurement}
\label{sec:method-fit}

The first rule is the one I broke most often, and the one that is easiest to break without noticing, because
fitting feels like measuring. You take a record, you choose a model, and the model returns a number with a
residual beside it; the residual looks like an error bar, the number looks like a result, and the fact that the
model \emph{was} the source of the number is invisible in the output.

It is not invisible in the logic. If a quantity is extracted by fitting a model to a record, then that quantity
is a property of the pair (record, model), not of the record alone. Change the model and the number changes; the
only defence is that the model be forced by something outside the record, and that its misfit be reported rather
than absorbed. A quantity recovered by fitting is therefore not an independent measurement of anything. It is a
projection of the record onto a chosen family of curves, and it can be perfectly repeatable, perfectly smooth in
the fit parameter, and wrong by a factor.

The clearest instance in this report is the attenuation of the microstrip line, where two different models of
the same record return rates that differ by a factor of $1.86$: the two-wave fit returns about
$0.0215$\,dB/mm against an audited $0.0116$\,dB/mm for the work the walls actually removed
(Section~\ref{sec:msres-alpha}). Both fits converge; neither reports a large residual; the fit is simply
describing a record that does not contain what the model says it contains. The audited figure is not a fit at
all -- it is a volume integral of $\sigma|\mathbf{E}|^2$ over the conductors, divided by the power entering the
section -- and that is why it is the one that is quoted, and why the fit is quoted only as a fit.

The same rule applies one level down, to quantities that are definitionally tied to the fitted ones. The
effective permittivity of Section~\ref{sec:tline-epseff} is $\eps_{\mathrm{eff}} \equiv (\beta c/\omega)^2$: once
$\beta$ is known, $\eps_{\mathrm{eff}}$ adds no information whatsoever and inherits the full uncertainty of the
fit that produced $\beta$. Reporting both is legitimate bookkeeping. Reporting both as \emph{two} numbers, or
worse, as agreement between two numbers, is a category error, because the second cannot disagree with the first.
This is why Section~\ref{sec:worked} lists $\eps_{\mathrm{eff}}$ as a restatement of the phase constant rather
than as a second result, and why the fit residual of that record, $0.0128$, is recorded as a property of the fit
and never used as an error bar on anything.

\subsection{Never use the quantity under test as its own oracle}
\label{sec:method-oracle}

A special case of the first rule deserves its own heading, because it survives review: an estimator can be
built out of the very assumption it is supposed to test. If you estimate transmission by assuming the field
envelope is constant, then your estimator has already asserted the result. It cannot detect a departure from a
constant envelope, because the departure is the thing the model was built to ignore. The estimator is not
biased in the statistical sense, where more data would eventually correct it; it is blind, and more data
sharpens the blindness.

This is what happened to the flux estimator that was for a while the report's headline measurement of loss. The
field was sampled at nine stations along the line, the time-averaged Poynting flux was formed at each, and the
square root of the last divided by the first was read as a transmission coefficient. On paper that is a
model-free measurement of attenuation. In fact the estimator assumes that the flux at a station is a smooth
function of position, declining because energy is being absorbed -- that is, it assumes what it is measuring.
Where the assumption holds it returns a plausible number; where it fails it returns a number of the same
plausible size, because the failure mode, a standing wave, is itself a smooth function of position. There is no
residual to warn you. Section~\ref{sec:worked} takes this record apart station by station and shows the reading
was dominated by a standing-wave ripple in the flux, not by absorption.

The same trap sits in every de-embedding step. Moving a reference plane by $d$ assumes a single travelling wave
of known $\beta$; if a reflected wave is also present, the move mixes them, and the mixing is smooth and
plausible. The rule that follows is blunt: before trusting a derived number, ask what the derivation assumed,
and then ask whether that assumption is precisely the proposition the number is supposed to establish. If it
is, the number is decoration.

\subsection{Stability is not accuracy}
\label{sec:method-stability}

A ratio is the most reassuring object in numerical work, because it cancels. Systematic errors that scale both
numerator and denominator leave the ratio untouched, and a ratio that moves by a part in a thousand between runs
feels like a measurement accurate to a part in a thousand. That feeling is wrong whenever the thing that is
stable is the \emph{scale} rather than the \emph{shape}.

Take two numbers that are both short by the same factor: $a \to \alpha a$ and $b \to \alpha b$. Their ratio
$a/b$ is unchanged, exactly, and remains unchanged as $\alpha$ wanders over any range at all. The ratio is
perfectly stable and completely silent about $\alpha$. Every quantity in this report that is extracted as a
ratio is exposed to this: the reference impedance from the ratio of the two modal voltages, the transmission
coefficient from the ratio of two fluxes, the reflection coefficient from the ratio of two sampled amplitudes.
If a common gain error, a common cell-volume convention, or a common windowing factor sits on both terms of the
ratio, then the ratio is reproducible to many digits while the absolute level is wrong by tens of per cent.

The way to tell the two apart is to break the commonality deliberately. If the numerator and denominator are
computed by genuinely different operations -- a volume integral of dissipation on one side, a flux integral over
a surface on the other -- then a common factor is unlikely and the ratio means something. Where both terms come
from the same field dump through the same windowing and the same normalisation, the ratio measures the
arithmetic and not the physics. That test is what separates the audited attenuation of
Section~\ref{sec:ms-accounting} from the station ratios of Section~\ref{sec:ms-estimators}: the audit divides
joules lost by joules entering, computed from different parts of the field, while the station ratio divides two
samples of the same quantity taken by the same operation. The first is a measurement that happens to be a
ratio; the second is a ratio that is not a measurement.

\subsection{A tolerance needs a pair of measurements}
\label{sec:method-tolerance}

A single number has no error bar. It has a value, and the value is exact in the sense that it was printed, and
nothing about it says how far it might sit from the truth. A target with no error bar therefore cannot be
failed: the comparison of one number against a bound is a comparison of two conventions, and the convention
wins or loses by fiat rather than by evidence.

This is not a philosophical point, it is the reason the acceptance criteria of this report are written the way
they are. When the criterion of Section~\ref{sec:ms-criterion} asks whether a measured rate of loss sits within
a stated band, it can only do so because there are two independent routes to the rate -- the audited joule
budget and the fitted decay -- whose \emph{difference} is a quantity with a meaning. The difference can be
large or small; that is evidence. One number compared against a target is not evidence of anything, and a
document full of such comparisons is a document that has never been tested.

The failure is common enough in this report's own history to be embarrassing. The reflection result of
Section~\ref{sec:msres-verdict} was for a while quoted as ``$-44.48$\,dB, goal $-20$\,dB met with $24$\,dB to
spare''. That is a true sentence and it is not a measurement, because the $24$\,dB is a statement about a
tolerance chosen in advance, not about the difference between two ways of obtaining the reflection. The moment a
second route was computed -- the summary line of the same dump carrying $|\gamma| = 0.00628$ against the
$0.00597$ of the plane rows, $5.2$ per cent apart -- the number acquired a real error bar, and the honest
statement became that the agreement is good to a few per cent and that the $5.2$ per cent itself is
unreconciled. Both statements describe the same run. Only the second one can be wrong.

The corollary is that a report which states tolerances must state them twice: once as the band it wishes to
meet, and once as the spread of the routes it actually has. Where there is only one route, the second number is
absent and the tolerance should be described as advisory, not met.

\subsection{Transients get a discard window, not a correction}
\label{sec:method-transients}

Every time-domain record begins before the physics does. The drive is switched on into a structure that is not
yet in its steady state, energy has not yet filled the section, and the reference plane sees a field that
contains the switch-on as well as the wave. There are two ways to deal with this, and only one of them is
honest.

The honest way is to discard. Accumulate the transform over a window that begins after the transient has left
the domain of interest, and report the start of the window as part of the provenance of the number; if the
window cannot be made long enough for the residual transient to be negligible, the number is provisional, not
corrected. The dishonest way is to model the transient -- to fit a decaying offset, subtract a mean, or
otherwise take the switch-on out of the record -- because the model of the transient is a hypothesis about the
same structure whose behaviour you are trying to measure. The correction and the measurement then share a
model, and the first rule applies again.

The record used throughout this report discards the first $600$ of $2400$ steps and accumulates over the
remaining $1800$, and that choice is stated with the number rather than buried in a script, because it is part
of what the number means. The window is not a convenience: a record that begins at step zero contains a
transient of the order of the resonator's filling time, and a transform over the whole record would mix it in.
Nothing about the arithmetic afterwards can remove it, so nothing is attempted; the window is the whole of the
treatment.

The same logic governs the drive. A resonant record is only as good as the frequency it was driven at, and the
report's records are driven at a snapped frequency that differs from the requested one by $-13.3$ parts per
million. That offset is small, and it is quoted anyway, because a resonance measured at a frequency that is not
the one you asked for is a resonance of a slightly different structure, and the reader is entitled to know how
different. Correcting the result to the requested frequency would require the very dispersion model that the
result is supposed to test.

\subsection{A repeatable diagnostic is still a diagnostic}
\label{sec:method-diagnostic}

The last rule is the one this report needed longest to learn, and it concerns a class of quantity that is
\emph{perfectly} well behaved. A diagnostic quantity is one that is computed by a clean operation from a clean
record, returns the same value every time the record is reprocessed, varies smoothly with the run parameters,
and is entirely useless for the purpose its name suggests. Its repeatability is what makes it dangerous: a
quantity that scattered would be caught immediately, whereas a quantity that is stable to five digits looks
like the best number in the document.

The station-profile tilt of Section~\ref{sec:worked} is such a quantity. It is the fractional change in the
time-averaged Poynting flux between the first and last station of a nine-point array, it is trivially
reproducible, and it was for a period quoted as a transmission coefficient and hence as a rate of loss. It is
not a rate of loss. The stations span about a third of a guide wavelength, so the standing wave that exists on any
real line -- including a lossless one -- is not averaged over the span at all; what the tilt measures is the
phase at which the array happens to be sampled. The reading is stable because the sampling is deterministic, not
because the physics is captured. Its correctness as an arithmetic average is not in question; its competence as
a measurement of attenuation is nil, and no amount of repetition can change that.

The rule that follows is a naming rule. A quantity may be reported as a diagnostic, with the word in its label
and the caveat in its paragraph, or it may be reported as a measurement. It may not be reported as a
measurement because it is stable, because it is convenient, or because it agrees with something else. And if a
later audit shows that a quantity which was published under a measurement name in fact measures a sampling
artefact, the correct response is to relabel it and keep the value -- the value is a fact about the record and
should not be destroyed -- while withdrawing every conclusion that leaned on the name. That is the procedure
applied to the tilt in Section~\ref{sec:worked}, and to the two ``field loss'' figures that preceded it, both of
which are withdrawn in favour of the audit that replaced them.

\subsection{What it cost, and what it bought}
\label{sec:method-cost}

The rules above are cheap to state and were expensive to learn. The price of getting them wrong is not an error
bar on one number; it is a whole chain of reasoning built on a noun. A derived quantity that is called a rate of
loss will be compared against an audited rate of loss, will agree with it to a couple of per cent, and the
agreement will be written down as independent confirmation -- when in fact the two are the same arithmetic on
the same field, and the agreement is a coincidence or an identity. Withdrawing the second as a witness leaves
the first standing, which is the correct outcome, but it also removes a paragraph of the argument, and the
paragraph has to be rewritten rather than deleted.

Section~\ref{sec:worked} is the ledger of that process for four numbers, with the status of each recorded as
quotable, diagnostic only, or withdrawn. It is deliberately unflattering. The report has a much smaller number
of real measurements than its length suggests, and after the audit they are: cavity frequencies against an
analytic spectrum, a reflection coefficient against a matched reference plane, an attenuation against a joule
audit, and a scattering width against a Mie series -- each of them a comparison of two independent routes to one
quantity, which is the only shape in which a number can be defended. Everything else in this document is either
a restatement of one of these, a diagnostic that is named as one, or arithmetic that would fail the tests above
if it were called a measurement. That is not a comfortable inventory for a document of this size. It is,
however, an inventory that can be checked, and a number that can be checked is worth ten that cannot.
